\begin{proofbox}

	\textbf{\textcolor{CProof}{证明方法一：由距离范围与三角不等式求极值}}\par\smallskip\textbf{思路提示}
	先把椭圆上点到圆心的距离范围确定为 $[d_{\min},d_{\max}]$；
	再利用圆周半径恒为 $r$，把椭圆到圆周的最小距离转化为
	\[
		\min_{t\in[d_{\min},d_{\max}]}|t-r|.
	\]
	最大距离则由三角不等式得到。

	\medskip

	\textbf{\textcolor{CProof}{正式推导}}

	\begin{description}[style=nextline, leftmargin=2.8em, labelsep=0.5em, itemsep=0.55em]

		\item[\textbf{步骤一（参数化表示）：}]
		      设椭圆参数点为
		      \[
			      P(a\cos\theta,b\sin\theta),\qquad \theta\in[0,2\pi].
		      \]
		      则
		      \[
			      PM^2=(a\cos\theta-x_0)^2+(b\sin\theta-y_0)^2.
		      \]
		      这是关于 $\theta$ 的连续函数，因此 $PM$ 也是关于 $\theta$ 的连续函数。

		      由于 $[0,2\pi]$ 是闭区间，所以 $PM$ 一定存在最大值和最小值，分别记为
		      \[
			      d_{\max}=\max_{P\in E}PM,\qquad d_{\min}=\min_{P\in E}PM.
		      \]

		      \par\textbf{理由：}
		      闭区间上的连续函数一定能取到最大值和最小值。这里必须使用 $[0,2\pi]$，不能写成 $[0,2\pi)$，因为 $[0,2\pi)$ 不是闭区间。

		\item[\textbf{步骤二（说明 $PM$ 的取值是整个区间）：}]
		      椭圆是一条连通闭曲线，函数
		      \[
			      P\mapsto PM
		      \]
		      连续，因此 $PM$ 的取值集合也是一个连续区间。

		      既然 $PM$ 的最小值为 $d_{\min}$，最大值为 $d_{\max}$，所以椭圆上点到圆心的距离取值范围正是
		      \[
			      [d_{\min},d_{\max}].
		      \]

		      \par\textbf{理由：}
		      连续函数在连通集合上的像仍然连通；在实数轴上的连通集合就是区间。因此 $PM$ 会取遍 $d_{\min}$ 与 $d_{\max}$ 之间的所有中间值。

		\item[\textbf{步骤三（固定 $P$，求到圆周的最近距离）：}]
		      固定椭圆上一点 $P$。对任意 $Q\in\Gamma$，都有
		      \[
			      QM=r.
		      \]
		      由反向三角不等式，
		      \[
			      PQ\ge |PM-QM|=|PM-r|.
		      \]
		      并且当 $P,M,Q$ 三点共线，且 $Q$ 取在合适方向的圆周点上时，可以取到等号。

		      所以对固定的 $P$，
		      \[
			      \min_{Q\in\Gamma}PQ=|PM-r|.
		      \]
		      因此
		      \[
			      D_{\min}
			      =\min_{P\in E}\min_{Q\in\Gamma}PQ
			      =\min_{P\in E}|PM-r|.
		      \]

		      \par\textbf{理由：}
		      圆周上点到圆心距离恒为 $r$，所以固定 $P$ 后，最近距离只由 $PM$ 与 $r$ 的差决定。

		\item[\textbf{步骤四（最小距离分段）：}]
		      由步骤二可知，$PM$ 的取值范围是
		      \[
			      [d_{\min},d_{\max}].
		      \]
		      于是
		      \[
			      D_{\min}=\min_{t\in[d_{\min},d_{\max}]}|t-r|.
		      \]

		      分三种情况讨论：

		      \[
			      \begin{cases}
				      r<d_{\min},                & \text{此时 } t-r>0,\ D_{\min}=d_{\min}-r, \\[1mm]
				      d_{\min}\le r\le d_{\max}, & \text{此时存在 } t=r,\ D_{\min}=0,          \\[1mm]
				      r>d_{\max},                & \text{此时 } r-t>0,\ D_{\min}=r-d_{\max}.
			      \end{cases}
		      \]

		      因此
		      \[
			      D_{\min}=
			      \begin{cases}
				      d_{\min}-r, & r<d_{\min},                \\[2mm]
				      0,          & d_{\min}\le r\le d_{\max}, \\[2mm]
				      r-d_{\max}, & r>d_{\max}.
			      \end{cases}
		      \]

		      \par\textbf{理由：}
		      绝对值函数 $|t-r|$ 表示数轴上 $t$ 到 $r$ 的距离，最小值就是点 $r$ 到区间 $[d_{\min},d_{\max}]$ 的距离。

		\item[\textbf{步骤五（最大距离推导）：}]
		      固定椭圆上一点 $P$。对任意 $Q\in\Gamma$，由三角不等式得
		      \[
			      PQ\le PM+QM=PM+r.
		      \]
		      当 $P,M,Q$ 三点共线，且 $M$ 位于 $P,Q$ 之间时，等号成立。

		      所以对固定的 $P$，
		      \[
			      \max_{Q\in\Gamma}PQ=PM+r.
		      \]
		      于是
		      \[
			      D_{\max}
			      =\max_{P\in E}\max_{Q\in\Gamma}PQ
			      =\max_{P\in E}(PM+r)
			      =d_{\max}+r.
		      \]

		      \par\textbf{理由：}
		      圆周上离 $P$ 最远的点在圆心 $M$ 的另一侧，此时距离为 $PM+r$。

	\end{description}


\bigskip

\textbf{\textcolor{CProof}{证明方法二：双参数三角代换（先消去圆周角，再对椭圆角求导）}}

本方法不预先使用三角不等式的几何结论，而是\textbf{从两点坐标和距离平方直接计算}。
注意椭圆与圆周上的点是各自独立运动的，必须使用\textbf{两个不同的参数} $\theta$、$\varphi$。

\begin{description}[style=nextline, leftmargin=2.8em, labelsep=0.5em, itemsep=0.55em]

    \item[\textbf{步骤一（分别对两条曲线设点）：}]
          设椭圆上一点及圆周上一点分别为
          \[
              P=(a\cos\theta,b\sin\theta),\qquad
              Q=(x_0+r\cos\varphi,y_0+r\sin\varphi),
          \]
          其中 $\theta,\varphi\in[0,2\pi]$，并定义
          \[
              U(\theta)=a\cos\theta-x_0,\qquad
              V(\theta)=b\sin\theta-y_0,
          \]
          \[
              F(\theta)=U(\theta)^2+V(\theta)^2.
          \]
          此处 $F$ 仅用于简化代数式，暂不对它作几何解释。

    \item[\textbf{步骤二（展开距离平方，分离圆周角）：}]
          根据两点间距离公式，
          \[
              \begin{aligned}
                  PQ^2={}&(U-r\cos\varphi)^2+(V-r\sin\varphi)^2\\
                  ={}&F(\theta)+r^2\\
                    &-2r\bigl(U\cos\varphi+V\sin\varphi\bigr).
              \end{aligned}
          \]
          这里 $U,V$ 都只依赖于 $\theta$；对于固定的 $\theta$，它们是常数，
          因而可\textbf{先对 $\varphi$ 求最值}。

    \item[\textbf{步骤三（用辅助角公式消去 $\varphi$）：}]
          由辅助角公式（或柯西不等式），
          \[
              -\sqrt{U^2+V^2}\le U\cos\varphi+V\sin\varphi
              \le\sqrt{U^2+V^2}.
          \]
          两端均可取到：若 $F(\theta)>0$，分别令
          \[
              (\cos\varphi,\sin\varphi)
              =\pm\frac{(U,V)}{\sqrt{F(\theta)}};
          \]
          若 $F(\theta)=0$，则括号中的式子恒为零，结论也成立。
          因为 $r>0$，所以固定 $\theta$ 后，
          \[
              \begin{aligned}
                  \min_{\varphi}PQ^2
                      &=F+r^2-2r\sqrt F=(\sqrt F-r)^2,\\
                  \max_{\varphi}PQ^2
                      &=F+r^2+2r\sqrt F=(\sqrt F+r)^2.
              \end{aligned}
          \]
          取非负平方根，得到
          \[
              \boxed{\min_{\varphi}PQ=\bigl|\sqrt{F(\theta)}-r\bigr|,\qquad
              \max_{\varphi}PQ=\sqrt{F(\theta)}+r.}
          \]
          \textbf{关键：两个角度的最值问题到此变成只含 $\theta$ 的问题。}

    \item[\textbf{步骤四（对剩下的 $F(\theta)$ 求导）：}]
          为求 $\sqrt{F(\theta)}$ 的取值范围，先求 $F(\theta)$ 的最值。
          因为平方根函数在 $[0,+\infty)$ 上单调递增，所以二者的最值顺序一致。
          直接求导：
          \[
              \begin{aligned}
                  F'(\theta)
                    ={}&-2a(a\cos\theta-x_0)\sin\theta\\
                       &+2b(b\sin\theta-y_0)\cos\theta\\
                    ={}&-2\bigl[(a^2-b^2)\sin\theta\cos\theta\\
                       &\phantom{-2\bigl[}{}-ax_0\sin\theta+by_0\cos\theta\bigr].
              \end{aligned}
          \]
          令 $F'(\theta)=0$，得到驻点方程
          \[
              \boxed{(a^2-b^2)\sin\theta\cos\theta
              -ax_0\sin\theta+by_0\cos\theta=0.}
          \]
          求出\textbf{全部}驻点，比较其 $F$ 值；由于 $F$ 可导且以 $2\pi$ 为周期，
          一整周的全局最大值和最小值一定会在这些驻点中取得。

    \item[\textbf{步骤五（一般位置可用万能代换求驻点）：}]
          当方程不能直接化简时，令
          \[
              t=\tan\frac{\theta}{2},\qquad
              \sin\theta=\frac{2t}{1+t^2},\quad
              \cos\theta=\frac{1-t^2}{1+t^2}.
          \]
          代入上面的驻点方程，再乘 $(1+t^2)^2$，整理得
          \[
              \boxed{\begin{aligned}
                  0={}&by_0t^4+2(a^2-b^2+ax_0)t^3\\
                     &+2(ax_0-a^2+b^2)t-by_0.
              \end{aligned}}
          \]
          将各实根还原为椭圆参数，并\textbf{单独检查 $\theta=\pi$}（此时
          $\tan(\theta/2)$ 无定义），比较全部候选位置的 $F$ 值。
          若 $x_0=0$ 或 $y_0=0$，通常优先化为关于正弦或余弦的二次式，不必强行解四次方程。
          最终记
          \[
              d_{\min}=\sqrt{\min_{\theta}F(\theta)},\qquad
              d_{\max}=\sqrt{\max_{\theta}F(\theta)}.
          \]

    \item[\textbf{步骤六（回到两条曲线，求全局距离极值）：}]
          因为 $\sqrt{F(\theta)}$ 连续，且 $\theta$ 走完一整周，
          所以它的取值恰好是区间 $[d_{\min},d_{\max}]$。
          步骤三的公式因而给出
          \[
              \begin{aligned}
                  D_{\min}
                      &=\min_{\theta}\bigl|\sqrt{F(\theta)}-r\bigr|
                       =\min_{s\in[d_{\min},d_{\max}]}|s-r|,\\
                  D_{\max}
                      &=\max_{\theta}\bigl(\sqrt{F(\theta)}+r\bigr)
                       =d_{\max}+r.
              \end{aligned}
          \]
          因此
          \[
              \boxed{D_{\min}=
              \begin{cases}
                  d_{\min}-r, & r<d_{\min},\\[1mm]
                  0, & d_{\min}\le r\le d_{\max},\\[1mm]
                  r-d_{\max}, & r>d_{\max},
              \end{cases}
              \qquad D_{\max}=d_{\max}+r.}
          \]
          \par\textbf{易错提醒：}
          当 $d_{\min}\le r\le d_{\max}$ 时，可能在 $F(\theta)=r^2$
          的位置取得 $D_{\min}=0$；\textbf{这些位置不一定满足 $F'(\theta)=0$}。
          因此正确顺序是\textbf{先用求导确定 $F$ 的整体值域，再求绝对值的最小值}，
          不能只把 $F$ 的驻点代入 $|\sqrt F-r|$ 就下结论。

\end{description}

	\medskip

	\textbf{\textcolor{CProof}{结论回扣}}
	两种证明均得到
	\[
		D_{\min}=
		\begin{cases}
			d_{\min}-r, & r<d_{\min},                \\[2mm]
			0,          & d_{\min}\le r\le d_{\max}, \\[2mm]
			r-d_{\max}, & r>d_{\max},
		\end{cases}
		\qquad
		D_{\max}=d_{\max}+r.
	\]

\end{proofbox}